\documentclass[11pt]{article}

\usepackage[a4paper,margin=27mm]{geometry}
\usepackage{amsmath,amssymb,amsthm,mathtools}
\usepackage{booktabs}
\usepackage{graphicx}
\usepackage[hidelinks]{hyperref}

\newtheorem{theorem}{Theorem}
\newtheorem{proposition}[theorem]{Proposition}
\newtheorem{corollary}[theorem]{Corollary}
\newtheorem{lemma}[theorem]{Lemma}
\newtheorem{assumption}[theorem]{Assumption}
\theoremstyle{definition}
\newtheorem{definition}[theorem]{Definition}
\theoremstyle{remark}
\newtheorem{remark}[theorem]{Remark}

\newcommand{\E}{\mathbb E}
\newcommand{\Pp}{\mathbb P}
\newcommand{\Var}{\operatorname{Var}}
\newcommand{\Cov}{\operatorname{Cov}}
\newcommand{\cL}{\mathcal L}
\newcommand{\cM}{\mathcal M}
\newcommand{\R}{\mathbb R}

\title{Accumulated-Variance Functionals of Hidden Markov and\\
Semi-Markov Processes: Finite-Sample Upper and Lower Bounds}
\author{Wilson Lombardo}
\date{}

\begin{document}
\maketitle

\begin{abstract}
We study two accumulated variance functionals for stationary two-state hidden
Markov or hidden semi-Markov process, observed through a single trajectory,
and separate identifying these observable quantities from recovering the
latent parameters behind them.  For bounded durations, inverting block
moments gives finite-sample upper bounds, worked out with fully explicit
constants for a capped-geometric Gaussian family; on fixed interior classes,
the corresponding two-point lower bounds have the same rate.  On prescribed
local persistence cells a long-lag estimator matches the regular and saturated
regimes---though only cell by cell, with a lag built around
the cell's own scale, not through any single estimator that adapts across
unknown scales.  Unbounded durations known only through two moments are a
separate matter: no uniform exponential bound is claimed there either.
\end{abstract}

\medskip
\noindent\textbf{Keywords:} accumulated variance; hidden Markov process;
hidden semi-Markov process; finite-sample concentration; minimax lower bound;
persistence.

\medskip
\noindent\textbf{2020 Mathematics Subject Classification:}
Primary 60K15; Secondary 60J10, 60E15, 62M05, 62M10.

\section{Introduction}

This paper concerns inference for two observable probability functionals from
one stationary trajectory: the variance accumulated over a finite prediction
horizon and its long-run per-unit-time limit.  The data length and prediction
horizon are kept separate throughout.  This distinction is essential because
estimation error is determined by the former, whereas the sensitivity of the
target may grow with the latter.

The first structural point is that identification of the target functional is
strictly weaker than identification of a latent parametrization.  We therefore
derive deterministic stability inequalities directly in observable moment
coordinates before propagating error through a latent inverse.  The resulting
bounds remain meaningful at some parameter configurations where full latent
recovery is unavailable, provided that the accumulated-variance functional is
still identified.

The two-state Markov model serves as a benchmark: its geometric covariance
structure exposes the finite-horizon sensitivity and permits a transparent
three-moment estimator.  The central semi-Markov result treats bounded
durations through residual-time Markovization and block-moment inversion.  In
the capped-geometric Gaussian family, the observable inverse, a uniform
identification margin, and a pseudo-spectral-gap lower bound are all calculated
explicitly rather than left as abstract regularity constants.

The comparison splits in two.  On fixed interior classes, a Bernstein bound
for observable moments and a two-point lower bound have the same powers of
data length, confidence level, and prediction horizon.  Along the symmetric
persistent path, a lag tied to the persistence time attains the corresponding
lower-bound powers on either side of saturation---but only cell by cell, with
the cell scale fixed in advance; we do not construct an estimator adaptive to
an unknown persistence scale.

The computational procedures are included to show that the theoretical
quantities are evaluable and to audit constants, indices, and saturation
thresholds.  They are not the basis of the contribution.  The paper also
records why two duration moments alone cannot support a uniform exponential
claim for unbounded sojourn laws.

\section{Observation experiment and estimands}

Let $S_t\in\{1,2\}$ be stationary and let the observations be conditionally
independent given the latent path, with state-dependent emission laws $Q_i$ and
\[
 \E(X_t\mid S_t=i)=\mu_i,
 \qquad \Var(X_t\mid S_t=i)=\sigma_i^2.
\]
Only $X_{1:T}=(X_1,\ldots,X_T)$ is observed.  To avoid a rate ambiguity, $T$
always denotes the data length and $N$ the horizon of the target
\[
 G_N(\theta):=\Var_\theta\!\left(\sum_{t=1}^N X_t\right),
 \qquad
 v_\infty(\theta):=\lim_{N\to\infty}N^{-1}G_N(\theta),
\]
whenever the latter limit exists.  The exact expression for $G_N(\theta)$ is
taken from Lemma~1 of the population analysis and is not rederived here.

For the two-state alternating semi-Markov model, complete sojourns
$D_i\sim F_i$ have probability mass $f_i$, mean $\tau_i$, and variance
$\nu_i^2$.  Consecutive sojourns are mutually independent.  The stationary
occupation probabilities and inclusive residual-life laws are
\begin{equation}
 \pi_i=\frac{\tau_i}{\tau_1+\tau_2},
 \qquad
 \Pp(R_i=r\mid S_1=i)=\frac{\Pp(D_i\ge r)}{\tau_i},\quad r\ge1.
 \label{eq:equilibrium}
\end{equation}
The relevant aperiodicity condition is the span-one condition on a complete
return cycle,
\begin{equation}
 \gcd\{d_1+d_2:f_1(d_1)f_2(d_2)>0\}=1.
 \label{eq:aperiodicity}
\end{equation}
Separate span-one assumptions on $F_1$ and $F_2$ are not equivalent to
\eqref{eq:aperiodicity} and are not, by themselves, sufficient.

\section{Identification}

\begin{definition}[Latent and functional identification]
Let $P_\theta^{(\ell)}=\cL_\theta(X_{1:\ell})$.  A latent parameter is
identified from $\ell$ observations, up to state permutation, when
\[
 P_\theta^{(\ell)}=P_{\theta'}^{(\ell)}
 \quad\Longrightarrow\quad
 \theta'=s\theta
\]
for a permutation $s$ of the two state labels.  A scalar functional $g$ is
identified when the weaker implication
$P_\theta^{(\ell)}=P_{\theta'}^{(\ell)}\Rightarrow
g(\theta)=g(\theta')$ holds.
\end{definition}

\begin{proposition}[The target is more identifiable than the latent model]
If $\E_\theta X_t^2<\infty$, then $G_N$ is identified by
$P_\theta^{(N)}$, irrespective of whether the latent parametrization is
identified.  If $\sum_{h\in\mathbb Z}|\Cov(X_0,X_h)|<\infty$, then
$v_\infty$ is identified by the full stationary observed law.  In general,
no fixed finite block identifies $v_\infty$ without structural or covariance
tail restrictions.
\end{proposition}

\begin{proof}
The first assertion follows from
\[
 G_N=\int\!\left(\sum_{t=1}^N x_t\right)^2dP^{(N)}
 -\left\{\int\!\sum_{t=1}^N x_t\,dP^{(N)}\right\}^2.
\]
The second follows from
$v_\infty=\gamma_0+2\sum_{h\ge1}\gamma_h$, where every fixed-lag
covariance is a moment of the observed law.  A finite block contains no
information about unrestricted covariances beyond its length.
\end{proof}

\subsection{Markov submodel}

Let
\[
 P=\begin{pmatrix}1-a&a\\ b&1-b\end{pmatrix},
 \qquad \lambda=1-a-b.
\]
For a known number of states, a full-rank transition matrix
($\lambda\ne0$) and linearly independent emission laws, the nonparametric
two-state HMM is identified from the law of three consecutive observations,
up to label permutation \cite{GassiatCleynenRobin2016}.  For two probability
measures, distinctness implies linear independence; hence $\mu_1\ne\mu_2$
is sufficient, although not necessary, because it implies $Q_1\ne Q_2$.
The convention $\mu_1<\mu_2$ fixes the label.

The assumptions are substantive.  If $Q_1=Q_2$, latent dynamics are
invisible.  If $\lambda=0$, the latent states are independent over time and
the observations reduce to an ordinary mixture; an unrestricted
nonparametric mixture decomposition need not be identified.  Conversely,
$\mu_1=\mu_2$ does not alone destroy identification if $Q_1\ne Q_2$, but it
does remove the state signal from ordinary autocovariances of $X_t$.

\subsection{Semi-Markov model}

We assume that the embedded chain alternates, so that successive sojourns are
in different states.  If embedded self-transitions are allowed, a single
calendar-time run in state $i$ can be split into an arbitrary number of
consecutive sojourns, and $F_i$ is not identified even if the complete latent
path $S$ were observed.

The Markov identification theorem cannot simply be transferred to an HSMM.
With duration bounds $D_i^{\max}<\infty$, the augmented state
$E_t=(S_t,R_t)$ is a finite Markov chain, but all augmented states sharing
the same value of $S_t$ also share the same emission law.  Thus the emission
family in the expanded HMM contains repeated measures and violates the
linear-independence hypothesis of the preceding theorem.  A dedicated
injectivity argument is required.

Moreover, an arbitrary infinite-support $F_i$ is not identified by any fixed
finite block.  For a block of length $\ell$, probability mass can be
redistributed among at least three duration points greater than $\ell$ while
preserving total tail mass and its first moment.  This leaves all latent block
probabilities of length $\ell$ unchanged, but can change the second duration
moment.  Consequently, from $P_\theta^{(\ell)}$ alone one cannot generally
identify $F_i$ or even $\nu_i^2$.  Autocovariances alone are weaker still and
need not separate $F_1$ from $F_2$.

For finite-sample HSMM inference we therefore impose one of the following
verifiable model restrictions:
\begin{enumerate}
 \item $F_i$ belongs to a fixed finite-dimensional family
       $F_i(\cdot;\eta_i)$ and a finite observed-block moment map is
injective up to labels; or
 \item $F_i$ has known finite support and the structured HSMM likelihood is
       injective up to labels; or
 \item a sieve $F_i^{(D)}$ is used and its truncation bias is explicitly
       retained.
\end{enumerate}
We write the quantitative local form of this requirement as
\begin{equation}
 \inf_{\vartheta\in U(\theta_0)}
 s_{\min}\{D M_H(\vartheta)\}\ge\kappa_{\rm id}>0,
 \label{eq:idmargin}
\end{equation}
where $M_H(\theta)$ is the selected vector of observable moments.  This is
an assumption to be verified for the chosen duration and emission families,
not a consequence of $\mu_1\ne\mu_2$ alone.

There is nevertheless a useful constructive sufficient condition.  Assume
that the two-component marginal mixture is identifiable within the selected
emission class:
\begin{equation}
 \pi_1G_1+\pi_2G_2=\pi_1'G_1'+\pi_2'G_2'
 \quad\Longrightarrow\quad
 (\pi_i',G_i')_{i=1}^2=(\pi_{s(i)},G_{s(i)})_{i=1}^2
 \label{eq:mixtureid}
\end{equation}
for some label permutation $s$.  This holds, for example, for many regular
finite-dimensional mixture families, but not for unrestricted probability
measures.

\begin{proposition}[Constructive HSMM identification under mixture
identifiability]
Assume \eqref{eq:mixtureid}, $G_1\ne G_2$, finite mean durations, stationary
initialization, and strict alternation.  Then the complete observed law
identifies $(G_1,G_2,F_1,F_2)$ up to label permutation.  If
$\operatorname{supp}(F_i)\subseteq\{1,\ldots,D\}$, the block law of length
$D+2$ is sufficient, though not necessarily minimal.
\end{proposition}

\begin{proof}
The marginal law identifies $(\pi_i,G_i)$ by \eqref{eq:mixtureid}.  Since two
distinct probability measures are linearly independent, their $m$-fold
tensor products are linearly independent; hence the mixture decomposition of
$\cL(X_{1:m})$ identifies every latent word probability
$\Pp(S_{1:m}=s)$.  Stationarity and alternation give
\[
 c:=\Pp(S_0=1,S_1=2)=\Pp(S_0=2,S_1=1)
   =\frac1{\tau_1+\tau_2},
 \qquad \tau_i=\frac{\pi_i}{c}.
\]
For $j\ne i$ and $d\ge1$,
\[
 f_i(d)=\frac{
 \Pp(S_0=j,S_1=\cdots=S_d=i,S_{d+1}=j)}{c}.
\]
All duration masses are therefore recovered from the complete law; a support
bound $D$ makes blocks through length $D+2$ sufficient.
\end{proof}

\begin{corollary}[Gaussian emission specialization]
If $G_i=\mathcal N(\mu_i,\sigma_i^2)$, $\pi_i\in(0,1)$, and
$(\mu_1,\sigma_1^2)\ne(\mu_2,\sigma_2^2)$, finite Gaussian-mixture
identifiability \cite{Teicher1963} verifies \eqref{eq:mixtureid}.  Hence the
complete observed law identifies
\[
 (\mu_i,\sigma_i^2,F_i,\tau_i,\nu_i^2)_{i=1}^2
\]
up to labels whenever the displayed duration moments exist.  Under
$\operatorname{supp}(F_i)\subseteq\{1,\ldots,D\}$, the block law through
length $D+2$ identifies the same tuple.  The restriction $\mu_1<\mu_2$
fixes labels when the means are separated; if the means coincide but the
variances differ, another deterministic labeling convention is required.
\end{corollary}

The separation condition $\mu_1\ne\mu_2$ alone is not a replacement for
\eqref{eq:mixtureid}.  At the i.i.d. latent boundary, different decompositions
of the same one-dimensional mixture can yield different emission and
geometric run-length parameters while producing the same i.i.d. observed
process.  Thus neither separated means nor nondegenerate sojourn laws suffice
on an unrestricted emission class.  If $G_1=G_2$, all latent dynamics are
invisible.

\section{Estimators}

\subsection{A model-free benchmark}

Let $\widehat\gamma_h$ be an empirical autocovariance.  For fixed $N<T$,
\begin{equation}
 \widehat G_N^{\rm dir}
 =N\widehat\gamma_0+2\sum_{h=1}^{N-1}(N-h)\widehat\gamma_h
 \label{eq:direct}
\end{equation}
estimates $G_N$ without identifying any latent parameter.  Deterministically,
\begin{equation}
 |\widehat G_N^{\rm dir}-G_N|
 \le N^2\max_{0\le h<N}|\widehat\gamma_h-\gamma_h|.
 \label{eq:directerror}
\end{equation}
This is an essential baseline, but its use of $N$ lags makes it inefficient
when $N$ grows with $T$.  For the long-run target we use the positive
semi-definite Bartlett HAC estimator of Newey and West
\cite{NeweyWest1987}, not the rectangularly truncated covariance sum.  With
$Y_t=X_t-\bar X_T$ and
\[
 \widetilde\gamma_h=T^{-1}\sum_{t=1}^{T-h}Y_tY_{t+h},
 \qquad
 w_{h,H}=1-\frac{h}{H+1},
\]
define
\begin{equation}
 \widehat v_H^{\rm B}
 =\widetilde\gamma_0+
 2\sum_{h=1}^{H}w_{h,H}\widetilde\gamma_h.
 \label{eq:bartlett-hac}
\end{equation}
Extending $Y_t$ by zero outside $\{1,\ldots,T\}$ gives the exact identity
\begin{equation}
 \widehat v_H^{\rm B}
 =\frac{1}{T(H+1)}
   \sum_{j=1}^{T+H}\left(\sum_{r=0}^{H}Y_{j-r}\right)^2\geq0.
 \label{eq:bartlett-psd}
\end{equation}
Thus the scalar estimator is nonnegative in every sample.  In contrast,
rectangular weights $w_h=1$ do not define a positive-semidefinite lag window
and can produce a negative estimate of a variance.  The stochastic error in
\eqref{eq:bartlett-hac} is of order $H$ times the maximal covariance error,
plus the covariance-tail and Bartlett-window biases.

Autocovariances alone do not generally identify the full HSMM tuple.  In
particular, $\gamma_0$ contains only the weighted combination
$\pi_1\sigma_1^2+\pi_2\sigma_2^2$, finitely many lags cannot identify an
infinite duration tail, and the scale $(\mu_1-\mu_2)^2$ can be confounded with
the latent indicator covariance.  A covariance-based GMM therefore requires
injectivity of its chosen finite-dimensional moment map; higher-order moments
or likelihood information may be necessary.

\subsection{A stable three-moment estimator for the Markov submodel}

Put
\[
 A=\pi_1\pi_2(\mu_1-\mu_2)^2,
 \qquad B=\pi_1\sigma_1^2+\pi_2\sigma_2^2.
\]
Then
\[
 \gamma_0=A+B,\qquad \gamma_h=A\lambda^h\quad(h\ge1).
\]
Fix $\rho<1$ and define
\[
 \mathcal Z_\rho=\{(A,B,\lambda):A,B\ge0,\ |\lambda|\le\rho\},
 \qquad M(A,B,\lambda)=(A+B,A\lambda,A\lambda^2).
\]
Rather than form the unstable ratio $\widehat\gamma_2/\widehat\gamma_1$,
let
\begin{equation}
 \widehat\zeta\in\arg\min_{\zeta\in\mathcal Z_\rho}
 \|M(\zeta)-(\widehat\gamma_0,\widehat\gamma_1,
 \widehat\gamma_2)\|_\infty,
 \qquad \widehat G_N^{\rm M}=G_N(\widehat\zeta).
 \label{eq:mdmarkov}
\end{equation}

\begin{theorem}[Functional stability without a signal lower bound]
\label{thm:functional-stability}
Suppose $\zeta,\zeta'\in\mathcal Z_\rho$ and set
$e_j=|M_j(\zeta')-M_j(\zeta)|$.  Define
\[
 S_{0,N}(\rho)=\sum_{k=0}^{N-2}(N-1-k)\rho^k,
\]
\[
 S_{1,N}(\rho)=\sum_{k=1}^{N-2}k(N-1-k)\rho^{k-1}.
\]
For $0\le \rho<1$ these coefficients contain no unresolved summation:
\begin{align}
 S_{0,N}(\rho)
 &=\frac{(N-1)-N\rho+\rho^N}{(1-\rho)^2},
 \label{eq:S0closed}\\
 S_{1,N}(\rho)
 &=\frac{(N-2)-N\rho+N\rho^{N-1}-(N-2)\rho^N}
         {(1-\rho)^3}.
 \label{eq:S1closed}
\end{align}
Then
\begin{align}
 |G_N(\zeta')-G_N(\zeta)|
 \le{}&Ne_0+2S_{0,N}(\rho)e_1
       +2S_{1,N}(\rho)(\rho e_1+e_2),
 \label{eq:stableGN}\\
 |v_\infty(\zeta')-v_\infty(\zeta)|
 \le{}&e_0+\frac{2e_1}{1-\rho}
       +\frac{2(\rho e_1+e_2)}{(1-\rho)^2}.
 \label{eq:stablevinf}
\end{align}
In particular, no lower bound on $A$ or $|\lambda|$ is required.
\end{theorem}

\begin{proof}
Write $u=A\lambda=\gamma_1$ and $v=A\lambda^2=\gamma_2$.  For $k\ge1$,
\[
 |u\lambda^k-u'{\lambda'}^k|
 \le \rho^k e_1+k\rho^{k-1}(\rho e_1+e_2),
\]
because
$|u'(\lambda-\lambda')|
\le\rho|u-u'|+|u\lambda-u'\lambda'|$.
Insert this inequality into
$G_N=N\gamma_0+2\sum_{k=0}^{N-2}(N-1-k)u\lambda^k$.
Letting $N\to\infty$ after normalization gives
\eqref{eq:stablevinf}; equivalently, sum the two geometric series.
\end{proof}

The theorem is deliberately functional-specific.  When $A=0$ or
$\lambda=0$, the latent decomposition is singular but every positive-lag
covariance vanishes, so the target remains stable.  This avoids an artificial
factor $1/|\gamma_1|$ introduced by ratio inversion.

\subsection{HSMM minimum distance and likelihood}

For a parametric or finite-support HSMM let
\[
 \widehat M_H=(\widehat m_1,\ldots,\widehat m_d),
 \qquad M_H(\theta)=\E_\theta\widehat M_H,
\]
where the moments may include marginal moments, autocovariances, and, when
needed for identification, third-order lagged moments.  Define
\begin{equation}
 \widehat\theta_{\rm MD}\in
 \arg\min_{\vartheta\in\Theta}
 \|W^{1/2}\{\widehat M_H-M_H(\vartheta)\}\|_2^2,
 \quad
 \widehat G_N=G_N(\widehat\theta_{\rm MD}),
 \quad
 \widehat v_\infty=v_\infty(\widehat\theta_{\rm MD}).
 \label{eq:mdhsmm}
\end{equation}
The exact $G_N(\widehat\theta)$ is evaluated by the coefficient-extraction or
renewal recursion of Lemma~1.

An alternative is the stationary HSMM maximum-likelihood estimator.  For
bounded durations, augment by residual time $E_t=(i,r)$.  Its transition
matrix has
\[
 Q_F((i,r),(i,r-1))=1\quad(r>1),
 \qquad
 Q_F((i,1),(3-i,d))=f_{3-i}(d),
\]
and stationary initial law
\begin{equation}
 \alpha_{i,r}=\Pp(S_1=i,R_i=r)
 =\frac{\Pp(D_i\ge r)}{\tau_1+\tau_2}.
 \label{eq:auginitial}
\end{equation}
Thus the first segment must be treated as left-censored through the
equilibrium residual law; the final segment is right-censored.  Standard
forward--backward recursions provide posterior state and segment counts.  The
Gaussian-emission M-step is a posterior-weighted update of $\mu_i$ and
$\sigma_i^2$, while duration probabilities are updated from posterior segment
counts with boundary censoring included in the likelihood.

For theory, $\widehat\theta$ should mean a global constrained likelihood or
minimum-distance minimizer.  EM is a numerical method for locating such a
point; without an initialization-basin argument it does not itself supply a
finite-sample guarantee.

\section{Non-asymptotic propagation}
\label{sec:nonasymptotic-propagation}

\begin{theorem}[Coordinatewise non-asymptotic delta method]
\label{thm:propagation}
Let $g:\Theta\to\R$ be continuously differentiable on a convex neighborhood
$U$ containing $\theta$ and $\widehat\theta$ on an event $E$.  Put
\[
 L_j=\sup_{\vartheta\in U}|\partial_jg(\vartheta)|.
\]
If
\[
 \Pp(|\widehat\theta_j-\theta_j|>t)
 \le\delta_j(t,T),
\]
then
\begin{equation}
 \Pp(|g(\widehat\theta)-g(\theta)|>\epsilon)
 \le \Pp(E^c)+
 \inf_{\substack{t_j>0\\\sum_jL_jt_j\le\epsilon}}
 \sum_j\delta_j(t_j,T).
 \label{eq:uniondelta}
\end{equation}
If, for all $x>0$, simultaneously with probability at least $1-2de^{-x}$,
\[
 |\widehat\theta_j-\theta_j|
 \le a_j\sqrt{x/T_{\rm eff}}+b_jx/T_{\rm eff},
\]
then on $E$
\begin{equation}
 |g(\widehat\theta)-g(\theta)|
 \le\sum_jL_j
 \left(a_j\sqrt{x/T_{\rm eff}}+b_jx/T_{\rm eff}\right).
 \label{eq:bernprop}
\end{equation}
\end{theorem}

\begin{proof}
On the intersection of the coordinate events, the multivariate mean-value
formula gives
$|g(\widehat\theta)-g(\theta)|\le\sum_jL_jt_j$.
The union bound proves \eqref{eq:uniondelta}; substituting the displayed
Bernstein radii proves \eqref{eq:bernprop}.
\end{proof}

\subsection{Observable-moment concentration and the Markov corollary}

The concentration step can be made explicit by Markovizing the semi-Markov
process.  Put $Z_t=(S_t,R_t,X_t)$ and, for a maximum lag $H$, put
$W_t^{(H)}=(Z_t,\ldots,Z_{t+H})$.  Under bounded duration support,
$Z=(Z_t)$ and each window process $W^{(H)}$ are stationary Markov chains.
Let $g_H>0$ denote the pseudo-spectral gap of $W^{(H)}$ (or any certified lower
bound for it).  This formulation is useful because the lag product is the
bounded additive functional
\[
 f_h(W_t^{(H)})=(X_t-m_X)(X_{t+h}-m_X),
 \qquad m_X=\E X_t.
\]
Thus concentration of empirical autocovariances is an application of an
existing Markov-chain additive-functional inequality, rather than a new
semi-Markov concentration principle.

The applications below use the corrected author version of
Paulin~\cite[Assumption~3.1 and Theorem~3.11, equation~(3.25)]{Paulin2015}.
Accordingly, each window chain is assumed time-homogeneous,
$\phi$-irreducible, and aperiodic on a Polish state space, with a unique
stationary law; the equilibrium initialization makes the observed window
stationary.  For a bounded measurable $f$ with $|f|\le K$, one has
$V_f:=\Var_\pi(f)\le K^2$ and $|f-\pi f|\le2K$.  Solving Paulin's quadratic exponent for
the empirical-average deviation gives, for $m\ge1$, $g>0$, $K>0$, and $y>0$,
\begin{equation}
 {\mathfrak b}_{\rm P}(m,g,K,y)
 =\frac{K}{mg}\left[
 20y+\sqrt{400y^2+8yg\{m+g^{-1}\}}
 \right],
 \label{eq:bradius}
\end{equation}
and hence
\[
 \Pp_\pi\!\left\{
 \left|m^{-1}\sum_{t=1}^m f(W_t)-\pi f\right|
 >{\mathfrak b}_{\rm P}(m,g,K,y)
 \right\}\le2e^{-y}.
\]
No unspecified Bernstein constant remains in this radius.

\begin{proposition}[Empirical autocovariances from a finite-support HSMM]
\label{prop:autocov-concentration}
Assume $|X_t-m_X|\le R$ almost surely and $T>H$.  Define the common-sample
autocovariances
\[
 \widetilde\gamma_h
 =\frac1{T-H}\sum_{t=1}^{T-H}
 (X_t-m_X)(X_{t+h}-m_X),\qquad 0\le h\le H,
\]
and let $\widehat\gamma_h$ replace $m_X$ by
$\overline X_T=T^{-1}\sum_{t=1}^T X_t$.  Set
\begin{align}
 b_\mu(x)&={\mathfrak b}_{\rm P}(T,g_0,R,x+\log4),\\
 b_{\gamma,H}(x)
 &={\mathfrak b}_{\rm P}(T-H,g_H,R^2,x+\log\{2(H+1)\}),\\
 r_{T,H}(x)&=b_{\gamma,H}(x)+2R b_\mu(x)+b_\mu(x)^2.
 \label{eq:moment-radius}
\end{align}
Then, for every $x>0$,
\begin{equation}
 \Pp\!\left\{
 \max_{0\le h\le H}|\widehat\gamma_h-\gamma_h|
 >r_{T,H}(x)\right\}\le4e^{-x}.
 \label{eq:momentbernstein}
\end{equation}
\end{proposition}

\begin{proof}
For each $h$, apply the chosen bounded additive-functional Bernstein
inequality to $f_h(W_t^{(H)})$ and take a union bound over $h=0,\ldots,H$.
This gives the term $b_{\gamma,H}(x)$.  Apply the same inequality to
$X_t-m_X$ to obtain
$|\overline X_T-m_X|\le b_\mu(x)$.  Finally, with
$d=\overline X_T-m_X$ and $Y_t=X_t-m_X$,
\[
 |(Y_t-d)(Y_{t+h}-d)-Y_tY_{t+h}|
 \le 2R|d|+d^2.
\]
A second union bound proves \eqref{eq:momentbernstein}.
\end{proof}

For Gaussian or other sub-Gaussian emissions, one may instead apply the
unbounded-function renewal inequality of
\cite{Adamczak2008,AdamczakBednorz2015} to the sub-exponential products, or
use clipped observations and add the analytically controlled clipping bias.
Finite emission variance by itself is not sufficient for
\eqref{eq:momentbernstein}.

\begin{corollary}[Direct finite-horizon HSMM variance bound]
\label{cor:direct-hsmm-bound}
Let $H=N-1<T$ and define $\widehat G_N^{\rm dir}$ by
\eqref{eq:direct}, using the common-sample autocovariances of
Proposition~\ref{prop:autocov-concentration}.  Then
\begin{equation}
 \Pp\!\left\{
 |\widehat G_N^{\rm dir}-G_N|>
 N^2r_{T,N-1}(x)\right\}\le4e^{-x}.
 \label{eq:direct-hsmm-probability}
\end{equation}
\end{corollary}

\begin{proof}
On the event in Proposition~\ref{prop:autocov-concentration},
\[
 |\widehat G_N^{\rm dir}-G_N|
 \le
 \left\{N+2\sum_{h=1}^{N-1}(N-h)\right\}r_{T,N-1}(x)
 =N^2r_{T,N-1}(x).
\]
\end{proof}

\begin{corollary}[Finite-sample Markov functional bound]
\label{cor:markov-bound}
Let $r_T(x)=r_{T,2}(x)$ denote the radius in
\eqref{eq:moment-radius}.
For the estimator \eqref{eq:mdmarkov}, with the convention that any minimizer
may be selected,
\begin{equation}
 \Pp\{ |\widehat G_N^{\rm M}-G_N|>
 2L_N^{\rm mom}(\rho)r_T(x)\}
 \le 4e^{-x},
 \label{eq:markovbound}
\end{equation}
where
\[
 L_N^{\rm mom}(\rho)=
 N+2S_{0,N}(\rho)+2(1+\rho)S_{1,N}(\rho).
\]
The analogous bound for $v_\infty$ replaces this constant by
\[
 L_\infty^{\rm mom}(\rho)=
 1+\frac2{1-\rho}+\frac{2(1+\rho)}{(1-\rho)^2}.
\]
\end{corollary}

\begin{proof}
The true moment vector is feasible, so minimum-distance optimality and the
triangle inequality give
$\|M(\widehat\zeta)-M(\zeta)\|_\infty\le2r_T(x)$.
Apply Theorem~\ref{thm:functional-stability} coordinatewise.
\end{proof}

Equivalently, the requested $\epsilon$--$\delta$ form is the computable bound
\begin{equation}
\begin{aligned}
 \Pp\{|\widehat G_N^{\rm M}-G_N|>\epsilon\}
 &\le \delta_{N,T}^{\rm M}(\epsilon),\\
 \delta_{N,T}^{\rm M}(\epsilon)
 &=
 \min\!\left[
 1,\,
 4\exp\!\left\{
 -\sup\left\{x>0:
 2L_N^{\rm mom}(\rho)r_{T,2}(x)\le\epsilon
 \right\}\right\}\right].
\end{aligned}
 \label{eq:markov-epsilon-delta}
\end{equation}
The right-hand side is one-dimensional and monotone, so it is evaluated
without simulation; throughout these inverse-radius formulas
$\sup\varnothing$ is defined as $0$, yielding the trivial bound $1$ when
$\epsilon$ is below the smallest certified radius.

For fixed $\rho<1$, $L_N^{\rm mom}(\rho)=O(N)$ and
$L_\infty^{\rm mom}(\rho)=O(1)$.  Consequently,
\[
 \widehat v_\infty-v_\infty=O_{\Pp}(T^{-1/2}),
 \qquad
 \widehat G_N-G_N=O_{\Pp}(N T^{-1/2}),
\]
under fixed mixing constants.  If $N=T$, the absolute error is therefore
$O_{\Pp}(\sqrt T)$, while the normalized error remains $O_{\Pp}(T^{-1/2})$.
An $O(T^{-1/2})$ claim for the \emph{absolute} error of $G_T$ would have the
wrong scale.

There are two distinct ill-conditioned boundaries.  As $|\lambda|\uparrow1$,
mixing slows, with an effective size of order $T(1-|\lambda|)$, and the
functional sensitivity contributes factors of order $(1-\lambda)^{-2}$ for
positive persistence.  A conservative moment argument gives
\[
 |\widehat v_\infty-v_\infty|
 =O_{\Pp}\{T^{-1/2}(1-\lambda)^{-5/2}\}.
\]
As $\lambda\to0$, by contrast, the full latent model approaches its
non-identifiable i.i.d. boundary.  Parameter recovery deteriorates, but the
functional-specific estimator above need not, because the serial covariance
itself vanishes.

\subsection{A concrete end-to-end HSMM theorem}

Fix a duration bound $D$ and a compact parameter class $\Theta_D$ of
stationary, strictly alternating, two-state HSMMs with
$\operatorname{supp}(F_i)\subseteq\{1,\ldots,D\}$.  Fix labels, for example by
$\mu_1<\mu_2$, and exclude zero duration masses required for irreducibility.
Let $\ell\ge D+2$ and choose a finite vector of bounded block features
\[
 \Phi=(\phi_1,\ldots,\phi_d):\R^\ell\longrightarrow[-1,1]^d,
 \qquad
 \Psi(\theta)=\E_\theta\{\Phi(X_{1:\ell})\}.
\]
From one trajectory, set $m=T-\ell+1$ and
\begin{equation}
 \widehat\Psi_T=\frac1m\sum_{t=1}^{m}\Phi(X_{t:t+\ell-1}),
 \qquad
 \widehat\theta_T\in\arg\min_{\vartheta\in\Theta_D}
 \|\widehat\Psi_T-\Psi(\vartheta)\|_\infty .
 \label{eq:blockmd}
\end{equation}
This is a fully specified global minimum-distance estimator; unlike EM, its
definition does not depend on a local numerical optimum.

For any estimand $g:\Theta_D\to\R$, define its observable identification
modulus
\begin{equation}
 \omega_{g,\Theta_D}(u)=
 \sup\{|g(\theta)-g(\vartheta)|:
 \theta,\vartheta\in\Theta_D,\ 
 \|\Psi(\theta)-\Psi(\vartheta)\|_\infty\le u\}.
 \label{eq:functionalmodulus}
\end{equation}
This quantity is finite by compactness and is numerically certifiable by
finite-dimensional global optimization.  It also records the correct
impossibility statement: if $\omega_{g,\Theta_D}(u)$ does not tend to zero as
$u\downarrow0$, then the functional is not uniformly identifiable from the
selected moments.

\begin{theorem}[End-to-end block-moment bound for a bounded-duration HSMM]
\label{thm:hsmm-end-to-end}
Assume that the stationary augmented window chain of length $\ell$ has
pseudo-spectral gap at least $\underline g_\ell>0$ uniformly on $\Theta_D$.
For $x>0$, let
\begin{equation}
 a_{T,\ell,d}(x)=
 {\mathfrak b}_{\rm P}\!\left(
 T-\ell+1,\underline g_\ell,1,x+\log(2d)\right).
 \label{eq:blockradius}
\end{equation}
Then, simultaneously for $g=G_N$ and $g=v_\infty$,
\begin{equation}
 \Pp_{\theta_0}\!\left[
 |g(\widehat\theta_T)-g(\theta_0)|
 >
 \omega_{g,\Theta_D}\{2a_{T,\ell,d}(x)\}
 \right]\le 2e^{-x}.
 \label{eq:hsmm-modulus-bound}
\end{equation}
If, whenever $\|\Psi(\theta)-\Psi(\vartheta)\|_\infty\le r_0$,
\begin{equation}
 \|\theta-\vartheta\|_2
 \le\kappa_{\rm id}^{-1}
 \|\Psi(\theta)-\Psi(\vartheta)\|_\infty
 \label{eq:global-inverse-margin}
\end{equation}
and $g$ is $L_g$-Lipschitz on $\Theta_D$, then, provided
$2a_{T,\ell,d}(x)\le r_0$,
\begin{equation}
 \Pp_{\theta_0}\!\left\{
 |g(\widehat\theta_T)-g(\theta_0)|
 >2L_g\kappa_{\rm id}^{-1}a_{T,\ell,d}(x)
 \right\}\le2e^{-x}.
 \label{eq:hsmmparamrate}
\end{equation}
\end{theorem}

\begin{proof}
Each coordinate of $\widehat\Psi_T$ is a bounded additive functional of the
length-$\ell$ augmented window chain.  The Bernstein inequality used in
\eqref{eq:bradius}, followed by a union bound over $d$ coordinates, gives
\[
 \|\widehat\Psi_T-\Psi(\theta_0)\|_\infty
 \le a_{T,\ell,d}(x)
\]
outside an event of probability at most $2e^{-x}$.  Since $\theta_0$ is
feasible in \eqref{eq:blockmd}, minimum-distance optimality and the triangle
inequality imply
\[
 \|\Psi(\widehat\theta_T)-\Psi(\theta_0)\|_\infty
 \le2a_{T,\ell,d}(x).
\]
The definition \eqref{eq:functionalmodulus} proves
\eqref{eq:hsmm-modulus-bound}.  Under
\eqref{eq:global-inverse-margin}, the preceding display gives
$\|\widehat\theta_T-\theta_0\|_2
\le2\kappa_{\rm id}^{-1}a_{T,\ell,d}(x)$; Lipschitz continuity of $g$
then proves \eqref{eq:hsmmparamrate}.
\end{proof}

In the same notation, Theorem~\ref{thm:hsmm-end-to-end} is exactly
\begin{equation}
 \Pp_{\theta_0}\{|g(\widehat\theta_T)-g(\theta_0)|>\epsilon\}
 \le\delta_{g,T}^{\rm HSMM}(\epsilon),
 \label{eq:hsmm-epsilon-delta}
\end{equation}
where
\[
 \delta_{g,T}^{\rm HSMM}(\epsilon)=
 \min\!\left[
 1,\,
 2\exp\!\left\{
 -\sup\left\{x>0:
 \omega_{g,\Theta_D}\{2a_{T,\ell,d}(x)\}\le\epsilon
 \right\}\right\}\right].
\]
Under \eqref{eq:global-inverse-margin}, one may replace the modulus inside
the supremum by
$2L_g\kappa_{\rm id}^{-1}a_{T,\ell,d}(x)$.

The theorem gives the requested deterministic observable-moment step as the
pathwise implication
\begin{equation}
 |g(\widehat\theta_T)-g(\theta_0)|
 \le
 \omega_{g,\Theta_D}\!\left(
 2\|\widehat\Psi_T-\Psi(\theta_0)\|_\infty\right),
 \label{eq:observable-deterministic}
\end{equation}
not merely as a proposed future calculation.

For Gaussian emissions a concrete bounded choice is a finite collection of
Fourier cylinder features
\[
 \phi_{u,c}(x_{1:\ell})=\cos(u^\top x_{1:\ell}),\qquad
 \phi_{u,s}(x_{1:\ell})=\sin(u^\top x_{1:\ell}).
\]
On a compact class with separated component parameters, the frequency grid
must be selected so that $\Psi$ is injective and
\eqref{eq:global-inverse-margin} is certified.  The latter can be audited by
a quantitative inverse-function argument using both a lower bound on the
smallest singular value of $D\Psi$ and a Lipschitz bound for $D\Psi$ inside
local neighborhoods, together with a global separation check outside those
neighborhoods.  A singular-value bound alone is not a global inverse theorem.
This audit is essential: mean separation alone does not imply a uniform
$\kappa_{\rm id}>0$.

For finite duration support, the augmented transition matrix should also
satisfy a quantitative geometric-ergodicity bound
\begin{equation}
 \|(Q_F-\Pi_F)^h\|\le C_{\rm mix}\rho_F^h,
 \qquad \rho_F<1,
 \label{eq:augmix}
\end{equation}
uniformly over the parameter class; it yields a computable lower bound for
$\underline g_\ell$.  Aperiodicity alone is qualitative and does not provide
a uniform constant in \eqref{eq:augmix}.

Archivo TeX - fragmento de identificación explícita

For reference, the previously derived alternating-renewal expression
\begin{equation}
 v_\infty=
 \frac{\tau_1\sigma_1^2+\tau_2\sigma_2^2}{T_c}
 +(\mu_1-\mu_2)^2
 \frac{\tau_2^2\nu_1^2+\tau_1^2\nu_2^2}{T_c^3},
 \qquad T_c=\tau_1+\tau_2,
 \label{eq:vinfhsmm}
\end{equation}
has fully explicit sensitivities.  With $\Delta=\mu_1-\mu_2$ and
$A_c=\tau_2^2\nu_1^2+\tau_1^2\nu_2^2$,
\begin{align*}
 \partial_{\mu_1}v_\infty&=2\Delta A_c/T_c^3,&
 \partial_{\mu_2}v_\infty&=-2\Delta A_c/T_c^3,\\
 \partial_{\sigma_i^2}v_\infty&=\pi_i,&
 \partial_{\nu_1^2}v_\infty&=\Delta^2\tau_2^2/T_c^3,\\
 \partial_{\nu_2^2}v_\infty&=\Delta^2\tau_1^2/T_c^3,&
 \partial_{\tau_1}v_\infty
 &=\frac{\tau_2(\sigma_1^2-\sigma_2^2)}{T_c^2}
 +\Delta^2\left(\frac{2\tau_1\nu_2^2}{T_c^3}
 -\frac{3A_c}{T_c^4}\right),\\
 \partial_{\tau_2}v_\infty
 &=\frac{\tau_1(\sigma_2^2-\sigma_1^2)}{T_c^2}
 +\Delta^2\left(\frac{2\tau_2\nu_1^2}{T_c^3}
 -\frac{3A_c}{T_c^4}\right).
\end{align*}
These derivatives may be inserted directly into
\eqref{eq:uniondelta}.

Archivo TeX - fragmento de cota superior compatible

Archivo TeX - fragmento de cota inferior de optimalidad

\input{numerical_validation_fragment_v2.tex}

Even with observed states, the number of complete cycles is approximately
$T/T_c$; duration parameters therefore have a baseline error scale
$\sqrt{T_c/T}$, before dimension and hidden-state separation are included.
The mean cycle length alone is not a mixing constant: residual-life tails,
minimum duration probabilities, and the gap of $Q_F$ also matter.

If $F_i$ is arbitrary with only finite mean and variance, an exponential
Bernstein bound cannot hold uniformly.  Exponential or bounded duration tails
are sufficient routes to geometric mixing.  Under polynomial tails one must
use regenerative/Fuk--Nagaev-type bounds, and estimating $\nu_i^2$ at a
root-number-of-cycles rate generally requires a fourth duration moment.
Likewise, finite emission variance alone is insufficient for exponential
concentration of empirical second moments.

\section{Prior-art map and scope of the contribution}

The following distinctions are important.
\begin{enumerate}
 \item Nonparametric HMM identification under full-rank transitions and
       linearly independent emissions is established in
       \cite{GassiatCleynenRobin2016}.  Spectral/moment algorithms with
       finite-sample guarantees are established for identifiable HMM classes
       in \cite{HsuKakadeZhang2012,AnandkumarHsuKakade2012}; nonparametric
       oracle/minimax inequalities appear in
       \cite{DeCastroGassiatLacour2016,AbrahamGassiatNaulet2023,AbrahamGassiatNaulet2025}.
 \item Bernstein, empirical-process, and regenerative inequalities for
       additive functionals of geometrically ergodic Markov chains are already
       available in
       \cite{Adamczak2008,AdamczakBednorz2015,BertailClemencon2010,Paulin2015}.
       The inequality actually used here is
       Paulin~\cite[Theorem~3.11, equation~(3.25), corrected author
       version]{Paulin2015}: it treats a stationary, not necessarily reversible
       chain satisfying Assumption~3.1 on a Polish state space, a bounded
       $L^2(\pi)$ summand, and a positive pseudo-spectral gap.  Its
       Proposition~3.4, equation~(3.9), converts uniform-ergodicity mixing time
       into the pseudo-spectral-gap lower bound used in the non-asymptotic and
       explicit capped-duration sections.  Since a
       bounded-duration HSMM becomes Markov after residual-time augmentation,
       and a lag product becomes additive after window augmentation,
       concentration of empirical autocovariances is covered at this generic
       level.
 \item Merlev\`ede, Peligrad and Rio
       \cite[Theorem~2]{MerlevedePeligradRio2009} give the corresponding
       strong-mixing antecedent for centered, uniformly bounded variables with
       $\alpha(k)\le e^{-2ck}$; stationarity is not required.  Their denominator
       contains the covariance proxy $mv^2$, a boundedness term, and a linear
       term multiplied by $(\log m)^2$.  This result is not used to justify the
       log-free Paulin radius above.
 \item Finite-sample error bounds for empirical autocovariance operators under
       strong mixing are available from Mollenhauer et
       al.\ \cite{MollenhauerEtAl2022}; their direct insertion into a HAC
       estimator is therefore background rather than a novelty claim.
 \item For observed Markov-chain functionals, non-asymptotic long-run-variance
       estimation predates this work: Gyori and Paulin
       \cite{GyoriPaulin2012} analyze a truncated autocovariance estimator.
       We therefore make no general first-estimator claim.
 \item Melnyk and Banerjee's spectral algorithm
       \cite{MelnykBanerjee2017} handles discrete HSMMs with bounded duration
       and full-rank conditions, but from a set of independent training
       sequences and with the aim of scoring a test sequence---a different
       statistical experiment and target from the present one, which starts
       from one trajectory and targets $G_N$ or $v_\infty$ directly.
 \item Barbu and Limnios \cite{BarbuLimnios2006} and Trevezas and Limnios
       \cite{TrevezasLimnios2010} prove consistency and asymptotic normality of
       HSMM maximum-likelihood estimators.  Farcomeni
       \cite{Farcomeni2025} derives exact score and information recursions.
       These are important computational and asymptotic antecedents, but not
       end-to-end high-probability functional bounds.
 \item Khorshidian et al.\ \cite{KhorshidianEtAl2016} prove strong consistency
       and asymptotic normality for moments of semi-Markov reward processes,
       but from an observed semi-Markov kernel---the hidden-state
       identification problem from emissions, which is our starting point,
       does not arise in their setting.
 \item Posterior concentration for semi-Markov kernels in Votsi et
       al.\ \cite{VotsiEtAl2021} concerns observed histories; hidden emissions
       and accumulated variance are outside that experiment.
 \item De Castro, Gassiat and Le Corff
       \cite{DeCastroGassiatLeCorff2017} study filtering and smoothing in a
       nonparametric \emph{HMM}; this paper is not an HSMM parameter-estimation
       theorem.
\end{enumerate}

We do not rely on a joint result by Douc, Moulines, and Rio, and no such
citation is retained.  The concentration step for the Markov component is
instead based on Paulin~\cite[Theorem~3.11]{Paulin2015}, with its state-space,
stationarity, boundedness, and pseudo-spectral-gap hypotheses stated in
Section~\ref{sec:nonasymptotic-propagation}.  Where the strong-mixing route is discussed, Merlev\`ede, Peligrad
and Rio~\cite[Theorem~2]{MerlevedePeligradRio2009} are cited only for that
statement.  Neither reference is used to claim a concentration theorem for
the present hidden semi-Markov functional; the functional bounds are proved
under the assumptions stated in their respective theorems.

Theorem~\ref{thm:hsmm-end-to-end} gives the generic composition for the stated
block-moment estimator.  Section~\ref{sec:explicit-capped-gaussian} makes this
composition quantitative:
for a named capped-geometric Gaussian family it gives an exact observable
inverse, a closed uniform $\underline\kappa_{\rm id,cg}$, a closed uniform
$\underline g_3$, and an end-to-end radius with Paulin's numerical constants.
The restriction to antipodal means and known common variance is explicit; the
result does not silently claim the four-free-emission-parameter model.

A matching lower bound and the sensitivity calculation for $G_N$ and
$v_\infty$ are given in Section~\ref{sec:minimax-lower-bound}.  On fixed
interior HMM classes, the root-data-length parameter rates are already covered
by Abraham, Gassiat and Naulet; propagation to the present smooth functionals
is not claimed as a new generic rate.  The upper and lower functional powers
match in $T$, $\log(1/\delta)$, and $N$ on the stated fixed interior box, but
the uniform Doeblin and global inverse constants are not claimed to be sharp.
On the symmetric persistent Markov path with fixed known emission
parameters, the long-lag estimator of
Section~\ref{sec:persistence-upper-bound} matches the two-point lower bound in
both the regular and saturated regimes on a local multiplicative cell.
The precise restrictions on this comparison are collected in
Section~\ref{subsec:optimality-scope}.

\section*{AI tool disclosure}

OpenAI Codex (GPT-5 family, accessed through the Codex desktop application,
most recently on 28 August 2026) assisted with mathematical exposition,
bibliographic discovery, code generation, and manuscript editing.  The author
retains full responsibility for the mathematical statements, citations,
computations, and accompanying code.

\begin{thebibliography}{99}

\bibitem{NeweyWest1987}
W.~K. Newey and K.~D. West (1987).
A simple, positive semi-definite, heteroskedasticity and autocorrelation
consistent covariance matrix.
\emph{Econometrica} \textbf{55}(3), 703--708.
\url{https://doi.org/10.2307/1913610}.

\bibitem{Teicher1963}
H.~Teicher (1963).
Identifiability of finite mixtures.
\emph{The Annals of Mathematical Statistics} \textbf{34}(4), 1265--1269.
\url{https://doi.org/10.1214/aoms/1177703862}.

\bibitem{GassiatCleynenRobin2016}
E.~Gassiat, A.~Cleynen and S.~Robin (2016).
Finite state space non parametric hidden Markov models are in general
identifiable.
\emph{Statistics and Computing} \textbf{26}, 61--71.
\url{https://doi.org/10.1007/s11222-014-9523-8}.

\bibitem{HsuKakadeZhang2012}
D.~Hsu, S.~M. Kakade and T.~Zhang (2012).
A spectral algorithm for learning hidden Markov models.
\emph{Journal of Computer and System Sciences} \textbf{78}, 1460--1480.
\url{https://doi.org/10.1016/j.jcss.2011.12.025}.

\bibitem{AnandkumarHsuKakade2012}
A.~Anandkumar, D.~Hsu and S.~M. Kakade (2012).
A method of moments for mixture models and hidden Markov models.
\emph{Proceedings of COLT}, PMLR \textbf{23}, 33.1--33.34.
\url{https://proceedings.mlr.press/v23/anandkumar12.html}.

\bibitem{DeCastroGassiatLacour2016}
Y.~De Castro, E.~Gassiat and C.~Lacour (2016).
Minimax adaptive estimation of nonparametric hidden Markov models.
\emph{Journal of Machine Learning Research} \textbf{17}(111), 1--43.
\url{https://jmlr.org/papers/v17/15-381.html}.

\bibitem{AbrahamGassiatNaulet2023}
K.~Abraham, E.~Gassiat and Z.~Naulet (2023).
Fundamental limits for learning hidden Markov model parameters.
\emph{IEEE Transactions on Information Theory} \textbf{69}(3), 1777--1794.
\url{https://doi.org/10.1109/TIT.2022.3213429}.

\bibitem{BretagnolleHuber1978}
J.~Bretagnolle and C.~Huber (1978).
Estimation des densit\'es: risque minimax.
\emph{S\'eminaire de probabilit\'es} \textbf{12}, 342--363.
\url{https://www.numdam.org/item/SPS_1978__12__342_0/}.

\bibitem{WolferKontorovich2019}
G.~Wolfer and A.~Kontorovich (2019).
Minimax learning of ergodic Markov chains.
\emph{Proceedings of Machine Learning Research} \textbf{98}, 904--930.
\url{https://proceedings.mlr.press/v98/wolfer19a.html}.

\bibitem{ThoppePrashanthBhat2024}
G.~Thoppe, L.~A. Prashanth and S.~P. Bhat (2024).
Risk estimation in a Markov cost process: lower and upper bounds.
\emph{Proceedings of Machine Learning Research} \textbf{235}, 48124--48138.
\url{https://proceedings.mlr.press/v235/thoppe24a.html}.

\bibitem{AbrahamGassiatNaulet2025}
K.~Abraham, E.~Gassiat and Z.~Naulet (2025).
Frontiers to the learning of nonparametric hidden Markov models.
\emph{Journal of Machine Learning Research} \textbf{26}(155), 1--75.
\url{https://www.jmlr.org/papers/v26/24-2230.html}.

\bibitem{Paulin2015}
D.~Paulin (2015).
Concentration inequalities for Markov chains by Marton couplings and
spectral methods.
\emph{Electronic Journal of Probability} \textbf{20}(79), 1--32.
\url{https://doi.org/10.1214/EJP.v20-4039}.
The general-state-space applications in this paper use the corrected author
version v5 at \url{https://arxiv.org/abs/1212.2015}.

\bibitem{Adamczak2008}
R.~Adamczak (2008).
A tail inequality for suprema of unbounded empirical processes with
applications to Markov chains.
\emph{Electronic Journal of Probability} \textbf{13}, 1000--1034.
\url{https://doi.org/10.1214/EJP.v13-521}.

\bibitem{AdamczakBednorz2015}
R.~Adamczak and W.~Bednorz (2015).
Exponential concentration inequalities for additive functionals of Markov
chains.
\emph{ESAIM: Probability and Statistics} \textbf{19}, 440--481.
\url{https://doi.org/10.1051/ps/2014032}.

\bibitem{BertailClemencon2010}
P.~Bertail and S.~Cl\'emen\c{c}on (2010).
Sharp bounds for the tails of functionals of Markov chains.
\emph{Theory of Probability and Its Applications} \textbf{54}, 505--515.
\url{https://doi.org/10.1137/S0040585X97984401}.

\bibitem{MerlevedePeligradRio2009}
F.~Merlev\`ede, M.~Peligrad and E.~Rio (2009).
Bernstein inequality and moderate deviations under strong mixing conditions.
In \emph{High Dimensional Probability V}, 273--292.
\url{https://doi.org/10.1214/09-IMSCOLL518}.

\bibitem{MollenhauerEtAl2022}
M.~Mollenhauer, S.~Klus, C.~Sch\"utte and P.~Koltai (2022).
Kernel autocovariance operators of stationary processes: estimation and
convergence.
\emph{Journal of Machine Learning Research} \textbf{23}(327), 1--34.
\url{https://www.jmlr.org/papers/v23/20-442.html}.

\bibitem{GyoriPaulin2012}
B.~M. Gyori and D.~Paulin (2012).
Non-asymptotic confidence intervals for MCMC in practice.
arXiv preprint; journal publication not verified.
\url{https://arxiv.org/abs/1212.2016}.

\bibitem{MelnykBanerjee2017}
I.~Melnyk and A.~Banerjee (2017).
A spectral algorithm for inference in hidden semi-Markov models.
\emph{Journal of Machine Learning Research} \textbf{18}(35), 1--39.
\url{https://www.jmlr.org/papers/v18/15-468.html}.

\bibitem{BarbuLimnios2006}
V.~Barbu and N.~Limnios (2006).
Maximum likelihood estimation for hidden semi-Markov models.
\emph{Comptes Rendus Math\'ematique} \textbf{342}, 201--205.
\url{https://doi.org/10.1016/j.crma.2005.12.013}.

\bibitem{TrevezasLimnios2010}
S.~Trevezas and N.~Limnios (2010).
Maximum likelihood estimation for general hidden semi-Markov processes with
backward recurrence time dependence.
\emph{Journal of Mathematical Sciences} \textbf{163}(3), 262--274.
\url{https://doi.org/10.1007/s10958-009-9675-9}.

\bibitem{Farcomeni2025}
A.~Farcomeni (2025).
Exact score and information matrix for panel hidden semi-Markov models.
\emph{Statistics and Computing} \textbf{35}, article 55.
\url{https://doi.org/10.1007/s11222-025-10585-y}.

\bibitem{KhorshidianEtAl2016}
K.~Khorshidian, F.~Negahdari and H.~A. Mardnifard (2016).
Estimation for discrete-time semi-Markov reward processes: analysis and
inference.
\emph{Methodology and Computing in Applied Probability} \textbf{18},
885--900.
\url{https://doi.org/10.1007/s11009-015-9468-1}.

\bibitem{VotsiEtAl2021}
I.~Votsi, G.~Gayraud, V.~S. Barbu and N.~Limnios (2021).
Hypotheses testing and posterior concentration rates for semi-Markov
processes.
\emph{Statistical Inference for Stochastic Processes} \textbf{24}, 707--732.
\url{https://doi.org/10.1007/s11203-021-09247-3}.

\bibitem{DeCastroGassiatLeCorff2017}
Y.~De Castro, E.~Gassiat and S.~Le Corff (2017).
Consistent estimation of the filtering and marginal smoothing distributions
in nonparametric hidden Markov models.
\emph{IEEE Transactions on Information Theory} \textbf{63}, 4758--4777.
\url{https://doi.org/10.1109/TIT.2017.2696959}.

\end{thebibliography}

\end{document}
